% Current arithmetic introduction, v18. Old ar17 labels retained for compatibility.
\section{序論}
\label{ar17:introduction}

一般型曲面の Lagrangian torus fibration を退化から構成するとき，
局所の作用周期が整数格子を定めることに加え，その格子を異なる成分の間で接続する必要がある．
本稿は，五つの平面へ退化する五次曲面に現れる標準 $K_5$ branch graph と，
五つの四直線 $K_4$ cores を固定し，この整数接着問題を扱う．
中心となる局所結果は，三本の高分母 rays を持つ正の整数 chart では，
正の行列成分の graph が必ず完全 matching になるという定理である．
この符号の決定を，node の有限な予算と反対側の corner の整除条件に組み合わせることで，
高分母 support を八つの五頂点 graph 型へ帰着する．
一方，局所には分母が無限に大きくなる正の整数族が存在する．
本稿では，この局所的な自由度と大域的な制約の関係を，
整数性，正値性，marked clutching，node 配分を保って調べる．

\subsection{整数データと二つの主結果}

cores を $V=\{0,1,2,3,4\}$ で標識し，共有 branch $ij$ の分母を
$q_{ij}=q_{ji}>0$，core $i$ の rays $j,k$ の間の整数係数を $p^i_{jk}$ とする．
局所の基本方程式は
\[
 \sum_{k\ne i,j}p^i_{jk}=-q_{ij},\qquad
 p^i_{jk}p^i_{kj}=q_{ij}q_{ik}-1.
\]
さらに正の実数 $w^i_j$ が存在し，
\[
 w^i_k-\frac{p^i_{jk}}{q_{ij}}w^i_j>0
\]
となることを要求する．この不等式は puncture basins の正面積に対応する．
共有 branch では
\[
 q_{ij}\mid p^i_{jk}+p^j_{ik},\qquad
 n_{ij;k}=1-\frac{p^i_{jk}+p^j_{ik}}{q_{ij}}\in\ZZ_{\ge0},
\]
\[
 D^i_j+D^j_i=0,\qquad
 D^i_j=-q_{ij}-\sum_{k\ne i,j}p^i_{kj}
\]
を課す．正確な定義は節~\ref{ar17:setup}に置く．
局所方程式は分母の奇数性を含意し，両端の row sums から各 branch の node 総数は五となる．
$q_{ij}>1$ の辺からなる graph を高分母 support $H$ と呼ぶ．

\begin{theorem}[三高分母の完全 matching 性]
\label{ar18:main-matching}
一つの core の分母が $(1,q_1,q_2,q_3)$，$q_1,q_2,q_3>1$ であり，
対角成分零の整数行列 $P=(p_{jk})$ が row sums，pair products，
および正の weights の条件を満たすとする．
$p_{jk}>0$ のとき無向辺 $jk$ を置く graph は，四頂点の完全 matching である．
従って，低 ray の row には正成分がちょうど一つ，負成分が二つ現れる．
\end{theorem}

定理~\ref{mh18:matching}でこの主張を証明する．
分母の上限，大域 clutching，node 条件，対向 $D$ 相殺は仮定しない．
まず整数 row/pair equations から，正符号 graph の候補を matching，二つの道型，
および低 ray を含む三角形型に限る．道型は残る正の entry に関する二次式の係数の符号で排除する．
三角形型では，strict cycle-gain inequalities と整数性を併用する．
積の差が一であることから，実不等式が整数 entry の片側 bound に変わり，
残る三つの pair equations は正整数の恒等式と両立しなくなる．
ここでは正の weights の仮定が実際に必要であり，row/pair equations だけを満たす
正三角形型の整数例は存在する（注意~\ref{mh18:positivity-essential}）．

\begin{theorem}[高分母 support の八型への帰着]
\label{ar17:main-support}
定義~\ref{ar17:admissible}の五次の算術的接着データを考える．
$H\ne\varnothing$ なら，$H$ は五頂点全てを使う連結 graph であり，
三角形を含み，少なくとも六辺を持ち，$b_1(H)\ge2$ となる．
高分母には少なくとも二種類の値が現れる．
さらに，$H$ の各二価頂点は三角形に属し，$|E(H)|=6$ なら $H$ は葉を持つ．
残る graph 必要候補は，標識を忘れて八型，$K_5$ の標識を保って二百四十六個である．
特に，一つの core に incident な四分母が全て一なら，全ての分母が一である．
\end{theorem}

この定理は，本稿の graph 障害を合わせた結論である
（系~\ref{sf18:remaining-support}）．
八型という数は整数接着データや幾何的 atlas の存在数ではなく，
証明した必要条件を通過する support の数を表す．
本稿の対象は指定した cores，標準 cones，marked punctures を保つ整数 atlas の接着であり，
一般型曲面上の任意の Lagrangian torus fibration の分類ではない．

\subsection{matching 性が大域の接着に与える制約}

三高分母 core の低 row の符号が定まることにより，その row を隣接する core の
node 数から復元する方法が有効になる．二高分母の標準形と共通低 puncture の比較も合わせ，
次の三型を排除する：
\[
 \begin{split}
 G_{5,5}&=\{01,02,03,12,14,34\},\\
 G_{5,10}&=\{01,02,03,12,13,24,34\},\\
 G_{5,12}&=K_4\text{ に一本の葉辺を加えた graph}.
 \end{split}
\]
$G_{5,5}$ は三角形と四角形が一辺を共有する house graph である
（定理~\ref{hs18:no-house}）．この型では二価 cores の比較と二本の低辺の $D$ 相殺だけで十分であり，
三高分母の完全 matching 定理そのものは必要としない．
$G_{5,10}$ の排除では，唯一の二価 core の標準形を，
共通低 puncture における通常接着と正負の例外接着に分ける
（定理~\ref{gten18:no-support}）．
通常型は一つの中央模型に強制されるが，その反対 corner の signed divisor が存在しない．
正例外型は $D$ 相殺と奇分母の parity に反する．
負例外型では，固定分母十三の無限共鳴族を全て残したまま，
node 予算と三本の辺の $D$ 相殺から矛盾を導く．
この無限族を全て残して扱うことが，局所整数模型の分類から大域の接着障害へ進む際の要点となる．

$G_{5,12}$ では，第五 core が一高分母の葉となるため，
全低分母 hub の $K_4$ 証明をそのまま適用することはできない．
一高分母標準形の自由度を保ち，隣接する三高分母 cores の matching 性と
共有 punctures の整数条件を用いて排除する
（定理~\ref{lh18:no-K4-leaf}）．
この結果は support 全体が $K_4$ と一本の葉辺からなる場合に関するものであり，
$K_4$ を部分 graph として含む全ての support を禁止する主張ではない．

これらに先立つ四使用頂点 $K_4$ の排除では，全低分母の第五 core を hub とし，
その node 条件と $D$ 相殺を通常辺上の正の流れへ書き換える
（定理~\ref{kg17:no-K4}）．
二例外は大きい負 reverse の位置に応じた場合分けで排除される．
四例外は，node 予算によって符号が安定する無限領域を有限個の型へまとめ，
残る六候補を判別式 $-47$ の局所二次式で排除する．
この $K_4$ 証明自体は，一般の三高分母 matching 性を使わずに成立する．
各証明で用いる局所方程式と大域条件を分けることで，
他の support に再利用できる部分を明確にする．

\subsection{幾何的な先行研究と本稿の範囲}

複素超曲面の高次元 pairs-of-pants による分解と，多面体的な底上の torus fibers は，
Mikhalkin\cite[Theorems~1 and~1$'$]{Mikhalkin}によって研究された．
Leung--Wan\cite{LeungWan}は一般型超曲面の退化に tropical geometry を用い，
正則最高次形式の Calabi--Yau 型開部分への集中と，対応する fibers の
漸近的な special Lagrangian 性を調べた．
Zhang\cite[Theorem~1.1]{Zhang}は，次数 $d\ge5$ の超曲面に微分同相な四次元多様体を，
$d(d-4)^2$ 個の四次元 pair-of-pants と K3 surfaces の部分集合に分解する結果を与えた．
これらは一般型における分解と fibration の幾何的な背景を与える．
本稿が固定する五つの四直線 cores の整数分母配置は，それらの分解に追加する問題である．

Evans--Mauri\cite{EM}は contact link と Liouville cone を通じて
局所 Lagrangian torus fibrations を構成し，条件付きの SNC divisor 近傍の模型を与えた．
本稿の四直線 core の解釈はこの方法を背景とする．
また Fernández de Bobadilla--Pełka\cite[Theorem~6.14]{FernandezBobadillaPelka}は，
与えられた Kähler form の exact modification を用いて A'Campo space の境界で
torus fibrations を構成し，symplectic connection によって近傍 fibers の
codimension-zero 部分へ運ぶ一般的方法を与えた．
maximal Calabi--Yau degenerations では，generic region の torus fibration と
体積充填・漸近的な幾何の結果を得ている
\cite[Theorems~1.1--1.2]{FernandezBobadillaPelka}．
本稿で問うのは，こうした幾何的構成と整合する特定の標準整数 cone charts の接着であり，
それぞれの存在定理だけから本稿の分母の剛性は従わない．
SNC 境界に沿う相対構成，作用切片を含む neck の接続，
曲面全体の fibration の存在は伴論文\cite{SurfaceGeometry}で扱う．

標準 $K_4$ cone を用いることは，その link の抽象 graph を指定することより強い．
整数 fan は一様 matroid $U_{3,4}$ の Bergman fan であり，
線形 tropicalization，matroid，hyperplane arrangement の関係は
Feichtner--Sturmfels\cite{FeichtnerSturmfels}，
Ardila--Klivans\cite{ArdilaKlivans}で体系的に扱われている．
本稿では，さらに face 上の作用微分から来る格子との同定を課す．
tropical atlas の重なりを ambient 整数アフィン写像の制限として定める条件は，
Shaw\cite[Definition~3.1]{ShawSurfaces}，
Jell--Rau--Shaw\cite[Definitions~2.1 and~5.1]{JellRauShaw}に明示されている．
この条件と，torus cycles の作用微分との適合性を同時に要求する点が，
本稿の整数 clutching の出発点である．

Cartwright\cite[Definitions~2.5 and~2.8]{CartwrightComplexes}の tropical complex は，
dual complex に交点数のデータを加え，局所交叉行列を用いて
線形関数と divisor theory を定める．
Gross--Siebert\cite{GrossSiebert}は，特異点を持つ実アフィン幾何と
追加の接着データから Calabi--Yau 退化を再構成する．
本稿の $p,q$ は実際の作用周期の格子から定まるので，
これらとの比較には交点数，アフィン層，特異集合，作用切片の対応を指定する必要がある．
本稿の整数条件が各理論の仮定を満たすことを示さずに，
それらの存在定理や有限性を分母へ適用することはしない．

\subsection{principal minors，cycle gains，整数性}

一つの core の行列を $A=\operatorname{diag}(q_j)+P$ と書けば，
\[
 A\mathbf1=0,\qquad \det A_{\{j,k\}}=1\quad(j\ne k)
\]
となる．非対称行列の principal-minor assignment と
principal minors 間の多項式関係は古典的な問題である．
Lin--Sturmfels\cite{LinSturmfels}は $4\times4$ 行列の関係式を求め，
同論文 Proposition~4 では集合分割による principal minors と cycle-sums の相互変換を与えた．
本稿の三角形の cycle products の消去や Schur 補行列も，
この一般的な代数的枠組みの中に位置付けられる．
本稿で追加して追跡するのは，固定した kernel vector $\mathbf1$，
個々の entry の整数性と整除性，strict basin inequalities，
および異なる五行列の marked punctures での整合である．

正の weights の存在判定は，既知の cycle-gain 条件の特殊化である．
$p_{jk}>0$ の有向辺に gain $p_{jk}/q_j$ を与えると，
$w_k>(p_{jk}/q_j)w_j$ が全て同時に成立することは，
全ての directed cycles の gain が一未満であることと同値になる．
max-times algebra では，これは maximum cycle geometric mean，
Kleene star，subeigenvector の理論で説明できる
\cite[\S2, Theorem~2.2 and Proposition~2.5]{SergeevSchneiderButkovic}．
strict な場合は，十分小さい $\varepsilon>0$ を取り，全ての gains を共通因子 $1+\varepsilon$ で拡大しても
cycle 条件を保てるため，通常の非 strict な subeigenvector 条件から回収できる．
本稿ではこの判定を初等的に示し，有理 weights と整数積による厳密検算に用いる．

一方，cycle 条件そのものは正符号 graph を matching に限定しない．
道上の gain inequalities や，十分小さい gains を持つ三角形には
正の weights が存在し得る．
定理~\ref{ar18:main-matching}の内容は，これらに
row sums と全ての pair products を整数として同時に課すと，
三高分母の場合には完全 matching だけが残るという点にある．
例えば三角形の排除では $HK=q_1q_2-1$ と strict な cycle inequality を合わせて
$(tS-H)K<1$ を得て，整数性から $H\ge tS$ と結論する．
この段階を実係数の一般的な principal-minor relation や
数値的な tolerance に置き換えることはできない．
既存の一般理論との関係をこのように区別し，本稿の各整数結果は明記した条件から直接証明する．

関連する整数行列の理論として，
Corrales--Valencia\cite[\S\S2--3]{CorralesValencia}の graph 上の arithmetical structures がある．
同理論は非正の非対角成分を持つ $Z$-matrices，$M$-matrices を用い，
固定した graph に対する有限性を扱う．
本稿の $A\mathbf1=0$ は形式上似た条件であるが，
高分母の局所行列には正の非対角成分が必要となる
（補題~\ref{lit10:nonpositive-q1}）．
従ってその有限性を，ここでの分母の一様上界として用いることはできない．

determinant 一の整数式は frieze patterns にも現れる．
Cuntz--Holm\cite[\S\S2--3 and~7]{CuntzHolmFriezes}は，
零を許す tame integral friezes の quiddity cycles の表示と，
離散集合上の零を持たない固定高さの friezes の有限性を与えた．
Cuntz--Holm--Pagano\cite{CuntzHolmPagano}は代数的数の環上での有限性を研究した．
frieze では隣接 minors の条件に加え，境界行，周期性，quiddity cycle の積条件を課す．
本稿の六つの face determinants からそれらの構造への写像は与えていないため，
同じ unimodular な式を持つことだけから分類や有限性を適用しない．

\subsection{局所の無限族，完全復元，残る問題}

局所には，任意の奇数 $q>1$ に対して分母 $(q,1,1,1)$ の正の chart が存在する
（系~\ref{q9:local-all-odd}）．
二高分母の全整数行列には signed-divisor 表示があり，
正値性は二本の高低 pairing の符号を決める
（命題~\ref{kl10:two-high-parametrization}，定理~\ref{ls11:two-high-signs}）．
固定した一方の奇分母に対する partner の集合は，七と十三の共鳴を除けば有限であるが，
この二つでは明示的な無限族が残る（定理~\ref{fp12:partner-structure}）．
従って，局所整数性から分母全体の有限性を結論することはできない．

三高分母についても，matching 性は全行列の有限性を意味しない．
matching 型から得られる Schur 三元一次系は特異になり得る．
実際，階数二の系から
\[
 (q_a,q_b,q_c,q_0)
 =(409+600n,\ 3493+5100n,\ 1303+1904n,\ 1),
 \qquad n\in\ZZ_{\ge0}
\]
という正の整数局所族が得られる（定理~\ref{sh14:singular-family}）．
matching 定理により，この Schur 表示と $\beta\le-1$ の結果
（定理~\ref{sm14:beta-negative}）は，今後，全ての正の三高分母模型へ適用できる．
しかし，$\beta\le-2$ の範囲で三つの高 $D$ が全て非負となり得るかという問題は未解決である．
matching 性から，追加の係数の符号原理まで証明されたと解釈してはならない．

別の有限性として，低 ray の行と列を両方固定すれば，
三高分母の全整数 completion は高々二つとなる
（命題~\ref{fc17:at-most-two}）．
同符号の low-row entries を選ぶと，二次係数が消えない一変数二次式が得られ，
その整数根と整除条件から全行列を復元できる．
同じ固定データに対する二つの正の標識付き completion も存在する
（注意~\ref{fc17:sharp}）．
ここでは低 column を動かす自由度を残していないため，上の無限族と矛盾しない．
また二つの標識付き解が，幾何学的に非同値な二つの charts を与えると主張するものではない．

大域の解析には素数ごとの合同条件も併用できる．
$\ell\mid q_{ij}$ の辺からなる素数 support $G_\ell$ は，
任意の奇素数 $\ell$ に対して $K_4$ と book $B_3$ を含まず，$G_3$ は matching となる
（定理~\ref{ps10:forbidden}）．
七辺以上の素数 support は $\ell=5$，$G_5=K_{2,2,1}$ の場合に限られ，
任意の七つの分母の gcd は五を割る
（定理~\ref{ps11:dense}，系~\ref{ps11:dense-gcd}）．
これらは高分母 support $H$ と異なる graph に対する制約であり，
その剰余解の存在は整数 lifting や atlas の存在を保証しない．

証明中の有限計算では，先に方程式から有限範囲または有限個の無限型代表を導く．
node 予算から得られる有限な場合分け，符号が安定する全無限領域，
および明示的な共鳴族の不等式による排除を区別して記述する．
全分母が百一以下なら全て一となる検算
（定理~\ref{ex10:bounded-exclusion}）は，指定した分母範囲についての結果である．
これを，本稿の上限を仮定しない graph 障害と混同しない．

中心的な未解決問題は，残る八型のいずれかが
異なる高分母を混在させた算術的接着データを持つかどうかである．
従って，一般の五次分母剛性，すなわち全ての分母が一に強制されるかという問題も未解決である．
今後は，全三高分母 core に使える matching 性と Schur 表示を，
一高・二高分母の標準形，素数 support，node 配分，固定低データの二次復元と組み合わせる．
整数候補が得られた場合には，局所の実面積と作用切片を同時に整合させる幾何的な課題が残る．
本稿の結果は，その実現可能性を調べるための局所符号分類と，
標準五次模型の support に対する厳密な必要条件を与える．
